\section{Experimental Setup}
\label{sec:setup}

We design experiments to answer three research questions:

\textbf{RQ1:} Is the structural efficiency measurement framework functionally correct?
Can it execute end-to-end on real model checkpoints and produce interpretable metrics?

\textbf{RQ2:} Do GRPO and DPO exhibit different Semantic Edit Proportions (SEP)
under KL-matched conditions? Does execution reward selectively concentrate
policy movement on semantic AST nodes?

\textbf{RQ3:} What confounds affect checkpoint-comparison studies of RL fine-tuning,
and how can they be detected and prevented?

These questions map directly to the Introduction's contributions: RQ1 validates
the framework; RQ2 tests the core hypothesis; RQ3 documents methodological lessons.

\subsection{Datasets}

We evaluate on two standard function-level code generation benchmarks:

\begin{table}[h]
\centering
\caption{Evaluation datasets}
\label{tab:datasets}
\begin{tabular}{llll}
\toprule
\textbf{Dataset} & \textbf{Problems} & \textbf{Language} & \textbf{Source} \\
\midrule
HumanEval+ & 164 & Python & evalplus \cite{Liu2023EvalPlus} \\
MBPP+ & 378 & Python & evalplus \cite{Liu2023EvalPlus} \\
\bottomrule
\end{tabular}
\end{table}

\textbf{Rationale:} HumanEval+ and MBPP+ are the canonical benchmarks used by all
referenced baselines (CodeRL, T\"{U}LU 3, DeepSeek-Coder), enabling
comparison with prior work. Their structured nature --- Python functions with
explicit input/output contracts and test suites --- makes AST analysis
well-defined. The evalplus harness provides reliable execution-based evaluation.
We use evalplus's stricter test augmentation (+ variants) throughout.

\subsection{Model}

\textbf{Base model:} DeepSeek-Coder-7B-instruct-v1.5 (\texttt{deepseek-ai/deepseek-coder-7b-instruct-v1.5},
HuggingFace). This model is a decoder-only code-specialized transformer at 7B
scale with permissive license, representing the state-of-the-art open-weight
code LLM family. Using the instruct-tuned variant as the starting point for
post-training is consistent with standard practice \cite{Lambert2024TULU}.

\subsection{Alignment Methods (Conditions)}

We compare three post-training conditions, all starting from the same base model:

\begin{table}[h]
\centering
\caption{Alignment method conditions}
\label{tab:conditions}
\begin{tabular}{lll}
\toprule
\textbf{Condition} & \textbf{Method} & \textbf{Reward Signal} \\
\midrule
GRPO-binary & TRL GRPOTrainer & +1.0 (pass) / 0.0 (fail) \\
GRPO-error-type & TRL GRPOTrainer & Error taxonomy reward \\
DPO & TRL DPOTrainer & Execution-oracle preference pairs \\
\bottomrule
\end{tabular}
\end{table}

Error taxonomy rewards: SyntaxError=0.1, RuntimeError=0.3, AssertionError=0.7, pass=1.0.
DPO uses passing solutions preferred over failing (intended design; see Limitation L4 in Section~\ref{sec:limitations} for actual implementation note).

\textbf{Control variables:} Identical base model, identical training problems
(CodeAlpaca/OSS-Instruct subset), identical compute budget (1000 training steps),
KL penalty matching (GRPO $\beta$=0.04, DPO $\beta$=0.1, tuned to produce
comparable KL trajectories), and identical evaluation harness.

\subsection{Execution}

\textbf{Sub-experiment h-e1 (Proof-of-Concept, RQ1):} A smoke test to verify the
measurement infrastructure. Executed on 6 HumanEval+ problems with synthetic
GRPO and DPO code completions (hand-crafted to have known structural properties),
confirming that AST edit distance, KL matching, and bootstrap CI execute correctly.

\textbf{Sub-experiment h-m1 (Mechanism Analysis, RQ2):} Full structural analysis
on KL-matched checkpoint pairs from GRPO and DPO training. Intended to use
27 matched pairs from 10+ distinct checkpoints. Due to checkpoint aliasing
(see Section~\ref{sec:aliasing}), effective sample size was $n_\text{eff} \approx 2$.
Results are reported as preliminary.

\subsection{Evaluation Metrics}

\textbf{Primary (framework diagnostic):}
\begin{itemize}
  \item \textbf{Semantic Edit Proportion (SEP):} Fraction of edits targeting CF+DF nodes (Section~\ref{sec:methodology})
  \item \textbf{Semantic edit distance:} ZSS edit distance restricted to CF+DF nodes
  \item \textbf{Structural efficiency:} Semantic edit distance per unit KL divergence
\end{itemize}

\textbf{Statistical tests:}
\begin{itemize}
  \item Mann-Whitney U test on SEP distributions (significance threshold $\alpha = 0.05$)
  \item Bootstrap 95\% CI on mean SEP differential (10,000 samples)
\end{itemize}

\textbf{Secondary (training confirmation):}
Pass@1 on HumanEval+ and MBPP+ --- used only to confirm training progressed,
not as the primary comparison metric.

\subsection{Compute Resources}

Experiments run on NVIDIA H100 NVL (95,830 MiB), single GPU per run.
Conda environment \texttt{youra-h-e1} and \texttt{youra-h-m1} (Python 3.10, PyTorch,
TRL v1.3.0, evalplus, zss, tree-sitter).
